% Options for packages loaded elsewhere
\PassOptionsToPackage{unicode}{hyperref}
\PassOptionsToPackage{hyphens}{url}
\documentclass[
]{article}
\usepackage{xcolor}
\usepackage[margin=1in]{geometry}
\usepackage{amsmath,amssymb}
\setcounter{secnumdepth}{-\maxdimen} % remove section numbering
\usepackage{iftex}
\ifPDFTeX
  \usepackage[T1]{fontenc}
  \usepackage[utf8]{inputenc}
  \usepackage{textcomp} % provide euro and other symbols
\else % if luatex or xetex
  \usepackage{unicode-math} % this also loads fontspec
  \defaultfontfeatures{Scale=MatchLowercase}
  \defaultfontfeatures[\rmfamily]{Ligatures=TeX,Scale=1}
\fi
\usepackage{lmodern}
\ifPDFTeX\else
  % xetex/luatex font selection
\fi
% Use upquote if available, for straight quotes in verbatim environments
\IfFileExists{upquote.sty}{\usepackage{upquote}}{}
\IfFileExists{microtype.sty}{% use microtype if available
  \usepackage[]{microtype}
  \UseMicrotypeSet[protrusion]{basicmath} % disable protrusion for tt fonts
}{}
\makeatletter
\@ifundefined{KOMAClassName}{% if non-KOMA class
  \IfFileExists{parskip.sty}{%
    \usepackage{parskip}
  }{% else
    \setlength{\parindent}{0pt}
    \setlength{\parskip}{6pt plus 2pt minus 1pt}}
}{% if KOMA class
  \KOMAoptions{parskip=half}}
\makeatother
\usepackage{color}
\usepackage{fancyvrb}
\newcommand{\VerbBar}{|}
\newcommand{\VERB}{\Verb[commandchars=\\\{\}]}
\DefineVerbatimEnvironment{Highlighting}{Verbatim}{commandchars=\\\{\}}
% Add ',fontsize=\small' for more characters per line
\usepackage{framed}
\definecolor{shadecolor}{RGB}{248,248,248}
\newenvironment{Shaded}{\begin{snugshade}}{\end{snugshade}}
\newcommand{\AlertTok}[1]{\textcolor[rgb]{0.94,0.16,0.16}{#1}}
\newcommand{\AnnotationTok}[1]{\textcolor[rgb]{0.56,0.35,0.01}{\textbf{\textit{#1}}}}
\newcommand{\AttributeTok}[1]{\textcolor[rgb]{0.13,0.29,0.53}{#1}}
\newcommand{\BaseNTok}[1]{\textcolor[rgb]{0.00,0.00,0.81}{#1}}
\newcommand{\BuiltInTok}[1]{#1}
\newcommand{\CharTok}[1]{\textcolor[rgb]{0.31,0.60,0.02}{#1}}
\newcommand{\CommentTok}[1]{\textcolor[rgb]{0.56,0.35,0.01}{\textit{#1}}}
\newcommand{\CommentVarTok}[1]{\textcolor[rgb]{0.56,0.35,0.01}{\textbf{\textit{#1}}}}
\newcommand{\ConstantTok}[1]{\textcolor[rgb]{0.56,0.35,0.01}{#1}}
\newcommand{\ControlFlowTok}[1]{\textcolor[rgb]{0.13,0.29,0.53}{\textbf{#1}}}
\newcommand{\DataTypeTok}[1]{\textcolor[rgb]{0.13,0.29,0.53}{#1}}
\newcommand{\DecValTok}[1]{\textcolor[rgb]{0.00,0.00,0.81}{#1}}
\newcommand{\DocumentationTok}[1]{\textcolor[rgb]{0.56,0.35,0.01}{\textbf{\textit{#1}}}}
\newcommand{\ErrorTok}[1]{\textcolor[rgb]{0.64,0.00,0.00}{\textbf{#1}}}
\newcommand{\ExtensionTok}[1]{#1}
\newcommand{\FloatTok}[1]{\textcolor[rgb]{0.00,0.00,0.81}{#1}}
\newcommand{\FunctionTok}[1]{\textcolor[rgb]{0.13,0.29,0.53}{\textbf{#1}}}
\newcommand{\ImportTok}[1]{#1}
\newcommand{\InformationTok}[1]{\textcolor[rgb]{0.56,0.35,0.01}{\textbf{\textit{#1}}}}
\newcommand{\KeywordTok}[1]{\textcolor[rgb]{0.13,0.29,0.53}{\textbf{#1}}}
\newcommand{\NormalTok}[1]{#1}
\newcommand{\OperatorTok}[1]{\textcolor[rgb]{0.81,0.36,0.00}{\textbf{#1}}}
\newcommand{\OtherTok}[1]{\textcolor[rgb]{0.56,0.35,0.01}{#1}}
\newcommand{\PreprocessorTok}[1]{\textcolor[rgb]{0.56,0.35,0.01}{\textit{#1}}}
\newcommand{\RegionMarkerTok}[1]{#1}
\newcommand{\SpecialCharTok}[1]{\textcolor[rgb]{0.81,0.36,0.00}{\textbf{#1}}}
\newcommand{\SpecialStringTok}[1]{\textcolor[rgb]{0.31,0.60,0.02}{#1}}
\newcommand{\StringTok}[1]{\textcolor[rgb]{0.31,0.60,0.02}{#1}}
\newcommand{\VariableTok}[1]{\textcolor[rgb]{0.00,0.00,0.00}{#1}}
\newcommand{\VerbatimStringTok}[1]{\textcolor[rgb]{0.31,0.60,0.02}{#1}}
\newcommand{\WarningTok}[1]{\textcolor[rgb]{0.56,0.35,0.01}{\textbf{\textit{#1}}}}
\usepackage{graphicx}
\makeatletter
\newsavebox\pandoc@box
\newcommand*\pandocbounded[1]{% scales image to fit in text height/width
  \sbox\pandoc@box{#1}%
  \Gscale@div\@tempa{\textheight}{\dimexpr\ht\pandoc@box+\dp\pandoc@box\relax}%
  \Gscale@div\@tempb{\linewidth}{\wd\pandoc@box}%
  \ifdim\@tempb\p@<\@tempa\p@\let\@tempa\@tempb\fi% select the smaller of both
  \ifdim\@tempa\p@<\p@\scalebox{\@tempa}{\usebox\pandoc@box}%
  \else\usebox{\pandoc@box}%
  \fi%
}
% Set default figure placement to htbp
\def\fps@figure{htbp}
\makeatother
\setlength{\emergencystretch}{3em} % prevent overfull lines
\providecommand{\tightlist}{%
  \setlength{\itemsep}{0pt}\setlength{\parskip}{0pt}}
\usepackage{bookmark}
\IfFileExists{xurl.sty}{\usepackage{xurl}}{} % add URL line breaks if available
\urlstyle{same}
\hypersetup{
  pdftitle={Divergences and Pairs Plots},
  hidelinks,
  pdfcreator={LaTeX via pandoc}}

\title{Divergences and Pairs Plots}
\author{}
\date{\vspace{-2.5em}2026-04-17}

\begin{document}
\maketitle

\subsection{Introduction}\label{introduction}

A divergent transition is HMC's most informative warning: it tells you
that the leapfrog integrator could not navigate a region of the
posterior accurately, and the chain therefore did \emph{not} explore
that region. Unlike a high \(\hat R\), which can sometimes be cured by
running longer chains, a divergent transition signals a structural
problem with the posterior geometry that no amount of additional
sampling will fix. Diagnosing where divergences occur --- typically
through pairs plots that colour divergent draws --- is the path to a
reparameterisation that lets the sampler do its job.

\subsection{Prerequisites}\label{prerequisites}

A working understanding of Hamiltonian Monte Carlo, leapfrog
integration, and the geometric origins of divergent transitions;
familiarity with \texttt{bayesplot::mcmc\_pairs()} for visualising joint
posteriors.

\subsection{Theory}\label{theory}

The leapfrog integrator approximates Hamiltonian dynamics with finite
step size \(\epsilon\). In regions of high curvature (typically funnels
in hierarchical models, where a small variance parameter shrinks all
dependent effects), even a moderate \(\epsilon\) produces a Hamiltonian
error large enough that the proposal is biased. Stan flags such
proposals as divergent and refuses to use them, so the chain effectively
avoids the problem region instead of exploring it. The classical remedy
is a \emph{non-centred} reparameterisation: replacing
\(\beta_i \sim \mathcal N(\mu, \sigma^2)\) with
\(\tilde z_i \sim \mathcal N(0, 1)\) and
\(\beta_i = \mu + \sigma \tilde z_i\) decouples the funnel geometry and
makes the posterior easier to sample.

\subsection{Assumptions}\label{assumptions}

HMC or NUTS sampling; the divergent-transition flag is specific to these
algorithms. Random-walk samplers do not produce divergences but suffer
different pathologies (slow mixing, mode-trapping) that need separate
diagnostics.

\subsection{R Implementation}\label{r-implementation}

\begin{Shaded}
\begin{Highlighting}[]
\FunctionTok{library}\NormalTok{(brms); }\FunctionTok{library}\NormalTok{(bayesplot)}

\FunctionTok{data}\NormalTok{(sleepstudy, }\AttributeTok{package =} \StringTok{"lme4"}\NormalTok{)}

\NormalTok{fit }\OtherTok{\textless{}{-}} \FunctionTok{brm}\NormalTok{(Reaction }\SpecialCharTok{\textasciitilde{}}\NormalTok{ Days }\SpecialCharTok{+}\NormalTok{ (Days }\SpecialCharTok{|}\NormalTok{ Subject), }\AttributeTok{data =}\NormalTok{ sleepstudy,}
           \AttributeTok{chains =} \DecValTok{4}\NormalTok{, }\AttributeTok{iter =} \DecValTok{2000}\NormalTok{, }\AttributeTok{refresh =} \DecValTok{0}\NormalTok{,}
           \AttributeTok{control =} \FunctionTok{list}\NormalTok{(}\AttributeTok{adapt\_delta =} \FloatTok{0.99}\NormalTok{))}

\CommentTok{\# Divergence diagnostics}
\NormalTok{rstan}\SpecialCharTok{::}\FunctionTok{check\_hmc\_diagnostics}\NormalTok{(fit}\SpecialCharTok{$}\NormalTok{fit)}

\CommentTok{\# Pairs plot of parameters with divergences coloured}
\NormalTok{np }\OtherTok{\textless{}{-}} \FunctionTok{nuts\_params}\NormalTok{(fit)}
\FunctionTok{mcmc\_pairs}\NormalTok{(}\FunctionTok{as.array}\NormalTok{(fit), }\AttributeTok{pars =} \FunctionTok{c}\NormalTok{(}\StringTok{"b\_Days"}\NormalTok{, }\StringTok{"sd\_Subject\_\_Days"}\NormalTok{),}
           \AttributeTok{np =}\NormalTok{ np, }\AttributeTok{np\_style =} \FunctionTok{pairs\_style\_np}\NormalTok{(}\AttributeTok{div\_color =} \StringTok{"red"}\NormalTok{))}
\end{Highlighting}
\end{Shaded}

\subsection{Output \& Results}\label{output-results}

\texttt{check\_hmc\_diagnostics()} reports the count of divergent
transitions, tree-depth saturation events, and the energy Bayesian
fraction of missing information (E-BFMI). \texttt{mcmc\_pairs()}
overlays the joint draws of two parameters with red points marking
divergences; clustering of red points along the boundary of the
parameter space pinpoints the geometry that broke.

\subsection{Interpretation}\label{interpretation}

A reporting sentence: ``An initial fit produced 12 divergent transitions
concentrated near small values of the random-slope standard deviation;
switching to a non-centred parameterisation and increasing
\texttt{adapt\_delta} to 0.99 eliminated all divergences and yielded
\(\hat R = 1.00\) for every parameter.'' Report the diagnostic and the
fix together; a published Bayesian analysis with unmentioned divergences
is not a credible analysis.

\subsection{Practical Tips}\label{practical-tips}

\begin{itemize}
\tightlist
\item
  Any divergence is worth investigating; even a single one signals a
  region of the posterior the sampler did not explore correctly.
\item
  The classic funnel --- variance parameter near zero with correlated
  dependent effects --- is the most common offender; non-centred
  parameterisation is almost always the right fix.
\item
  Increasing \texttt{adapt\_delta} to 0.95 or 0.99 reduces step size and
  eliminates marginal divergences, but it is a band-aid: if the geometry
  is genuinely difficult, reparameterise.
\item
  \texttt{bayesplot::mcmc\_parcoord()} complements pairs plots by
  showing divergent transitions across many parameters simultaneously,
  which is helpful in models with dozens of variance components.
\item
  When divergences persist after non-centred parameterisation and high
  \texttt{adapt\_delta}, examine the prior --- improperly tight priors
  can pinch the posterior into geometry that no sampler can navigate.
\end{itemize}

\subsection{R Packages Used}\label{r-packages-used}

\texttt{bayesplot} for \texttt{mcmc\_pairs()},
\texttt{mcmc\_parcoord()}, and \texttt{nuts\_params()}; \texttt{rstan}
and \texttt{cmdstanr} for the underlying diagnostic outputs
(\texttt{check\_hmc\_diagnostics()}, \texttt{get\_sampler\_params()});
\texttt{posterior} for tidy access to NUTS metadata when building custom
diagnostic summaries.

\subsection{For Reviewers}\label{for-reviewers}

\textbf{What to look for in a paper using this method.}

\begin{itemize}
\tightlist
\item
  \textbf{Common misapplications.}
\item
  \textbf{Diagnostics that should be reported but often aren't.}
\item
  \textbf{Red flags in tables and figures.}
\item
  \textbf{What to verify.}
\item
  \textbf{What an adequate Methods paragraph must contain.}
\end{itemize}

\subsection{See also --- labs in this
chapter}\label{see-also-labs-in-this-chapter}

\begin{itemize}
\tightlist
\item
  \href{labs/lab-bayesian-thinking-control-vs-decision-error.Rmd}{Bayesian
  thinking; α control vs decision error}
\item
  \href{labs/lab-brms-stan-loo-hierarchical-models.Rmd}{brms / Stan,
  LOO, hierarchical models}
\end{itemize}

\end{document}
